\documentclass[10pt]{article}
\usepackage[T1]{fontenc}
\usepackage[utf8]{inputenc}
\usepackage[french]{babel}
\usepackage{amsmath,amssymb,mathrsfs,mathpazo}
\usepackage{geometry}
\geometry{a4paper,margin=2cm}
\usepackage{graphicx,array,tabularx,tabularray,colortbl,multirow,multicol,booktabs,enumitem,pifont,ulem,fancybox,diagbox,cancel}
\usepackage{pstricks,pstricks-add,pst-plot,pst-tree,pst-node,pst-text,pst-func,pst-math,pst-bspline,pst-3dplot,pst-eucl,pst-solides3d}
\let\pscalculate\undefined
\usepackage{pst-bezier}
\usepackage{tikz,pgfplots,tkz-tab}
\usetikzlibrary{arrows,arrows.meta,patterns,calc,positioning,trees,decorations.markings,shapes.geometric,shapes.misc,shapes.arrows,intersections,decorations.pathmorphing,decorations.pathreplacing,backgrounds,shadows,babel}
\pgfplotsset{compat=1.18}
\usepgfplotslibrary{fillbetween}
\usepackage[np]{numprint}
\usepackage{listings,makecell,pas-tableur,twoopt}
\let\barre\undefined
\usepackage[pstricks]{tablvar}
\definecolor{darkgreen}{RGB}{0,100,0}
\providecommand{\red}{\color{red}}
\providecommand{\blue}{\color{blue}}
\providecommand{\green}{\color{green}}
\providecommand{\black}{\color{black}}
\providecommand{\white}{\color{white}}
\providecommand{\gray}{\color{gray}}
\providecommand{\cyan}{\color{cyan}}
\providecommand{\magenta}{\color{magenta}}
\providecommand{\yellow}{\color{yellow}}
\providecommand{\violet}{\color{violet}}
\providecommand{\orange}{\color{orange}}
\providecommand{\euro}{\texteuro}
\providecommand{\EUR}{\texteuro}
\providecommand{\ieme}{\textsuperscript{e}}
\providecommand{\degres}{\ensuremath{^\circ}}
\providecommand{\degree}{\ensuremath{^\circ}}
\providecommand{\vv}[1]{\overrightarrow{#1}}
\providecommand{\Coord}[1]{#1}
\providecommand{\decosix}{}
\providecommand{\starredbullet}{\ensuremath{\bullet}}
\pagestyle{empty}
\setlength{\parindent}{0pt}
\usepackage{adjustbox}\DeclareUnicodeCharacter{2236}{:}\DeclareUnicodeCharacter{2190}{\ensuremath{\leftarrow}}
\providecommand{\R}{\mathbb{R}}
\providecommand{\N}{\mathbb{N}}
\providecommand{\Z}{\mathbb{Z}}
\providecommand{\Q}{\mathbb{Q}}
\providecommand{\D}{\mathbb{D}}
\providecommand{\C}{\mathbb{C}}
\providecommand{\e}{{\mathrm{e}}}
\providecommand{\vect}[1]{\overrightarrow{#1}}
\providecommand{\vectt}[1]{\overrightarrow{\mathrm{#1}}}
\providecommand{\vvt}[1]{\overrightarrow{\mathrm{#1}}}
\providecommand{\np}[2][]{{#2~#1}}
\providecommand{\vv}[1]{{\overrightarrow{#1}}}
\renewcommand{\d}{{\mathrm{d}}}
\providecommand{\degres}{{^\circ}}
\providecommand{\degre}{{^\circ}}
\providecommand{\blue}{\color{blue}}
\providecommand{\red}{\color{red}}
\providecommand{\enskip}{\quad}
\providecommand{\llbracket}{\lbrack\!\lbrack}
\providecommand{\rrbracket}{\rbrack\!\rbrack}
\providecommand{\euro}{{\text{€}}}
\providecommand{\no}{{\text{n°}}}
\providecommand{\interval}[2]{{[#1;#2]}}
\providecommand{\intervFO}[2]{{[#1;#2[}}
\providecommand{\dotfill}{\ldots}

\providecommand{\ds}{\displaystyle}
\providecommand{\textding}[1]{#1}
\providecommand{\Oij}{\(\mathrm{O};\vec{\imath},\vec{\jmath}\)}
\providecommand{\Oijk}{\(\mathrm{O};\vec{\imath},\vec{\jmath},\vec{k}\)}
\providecommand{\Ouv}{\(\mathrm{O};\vec{u},\vec{v}\)}
\providecommand{\intervFF}[2]{\ensuremath{[#1;#2]}}
\renewcommand{\euro}{\ifmmode\text{\texteuro}\else\texteuro\fi}
\DeclareUnicodeCharacter{00B0}{\ifmmode^{\circ}\else\textdegree\fi}
\DeclareUnicodeCharacter{20AC}{\ifmmode\text{\texteuro}\else\texteuro\fi}
% Énoncés : Éducation nationale. Transcription et corrigés : APMEP, Denis Vergès et contributeurs. Découpage et classement : CoopMaths, Sébastien Lozano, Éric Elter et contributeurs. Adaptation de la présentation : Maths Explorer.
% CC BY-SA : https://coopmaths.fr/www/about/
% https://coopmaths.fr/alea/static/bac/2022/tex/bac_2022_05_sujet2_madagascar_4_cor.tex
\begin{document}
\noindent\textit{APMEP / CoopMaths. CC BY-SA. Présentation adaptée par Maths Explorer.}\par\medskip
\medskip

\textbf{Partie A}

\medskip

\begin{enumerate}
	\item La fonction $x\mapsto \e^{-x}$ est continue en 0 et $\e^{-0}=1$. $\begin{gathered}\displaystyle\lim\limits_{\substack{x \rightarrow 0 \\ x>0}}\ln(x)=-\infty\end{gathered}$ donc $\begin{gathered}\displaystyle\lim\limits_{\substack{x \rightarrow 0\\ x>0}} f(x)=-\infty\end{gathered}$ 
	\item La fonction $f$ est continue et dérivable sur $]0~;~1]$. 
	
$\forall x \in ]0~;~1]$, $f'(x)=-\e^{-x}+\dfrac1x=\dfrac{-x\e^{-x}+1}{x}=\dfrac{1-x\e^{-x}}{x}$
	\item La fonction $x\mapsto \e^{-x}$ est continue, dérivable de dérivée $x\mapsto -\e^{-x}$. Or $\forall x \in ]0~;~1]$, $-\e^{-x}<0$. donc la fonction  $x\mapsto \e^{-x}$ est strictement décroissante sur $]0~;~1]$.
	
Donc $\forall x \in ]0~;~1]$, $:e^{-1} \leqslant \e^{-x}<\e^{0}<1$ donc $0< \e^{-x}<\e^{0}<1$. De plus $0<x\leqslant 1$ donc $0<x\e^{-x}<1$.
	
Cela signifie que $\forall x \in ]0~;~1]$, $1-x\e^{-x}>0$ donc $f'(x)>0$ donc $f$ est strictement croissante sur $]0~;~1]$. $f(1)=\e^{-1}+\ln(1)=\e^{-1}$. 
	
Ci-dessous le tableau de variation de la fonction $f$ :
	\begin{center}





\psset{unit=1cm}
\begin{pspicture}(6,3)
\psframe(6,3)\psline(0,2)(6,2)\psline(0,2.5)(6,2.5)\psline(2,0)(2,3)
\psline(2.1,0)(2.1,2.5)\psline(2.15,0)(2.15,2.5)
\uput[u](1,2.4){$x$}\uput[u](2.125,2.4){$0$}\uput[u](5.8,2.4){$1$}
\rput(1,2.25){$f'(x)$}\rput(4,2.25){$+$}
\rput(1,1){$f$}\uput[u](2.5,0){$- \infty$}\uput[d](5.6,2){$\text{e}^{- 1}$}
\psline{->}(2.6,0.5)(5.5,1.45)
\end{pspicture}
	\end{center}
\item La fonction $f$ est continue et strictement croissante sur $]0~;~1]$ à valeurs dans $]-\infty~;~\e^{-1}]$. Or $0 \in ]-\infty~;~\e^{-1}]$ donc d'après le corollaire du théorème des valeurs intermédiaires, l'équation $f(x)=0$ admet une unique solution,  notée $\ell$, dans $]0~;~1]$.

À la calculatrice : $\ell \approx 0,567$.
\end{enumerate}

\medskip

\textbf{Partie B}

\medskip

On définit deux suites $\left(a_n\right)$ et $\left(b_n\right)$ par:

\[\left\{\begin{array}{l c l}
	a_0& =& \dfrac{1}{10}\\
	b_0& =& 1
\end{array}\right.\: \text{et, pour tout entier naturel }\: n,\left\{\begin{array}{l c l}
	a_{n+1}&=&\text{e}^{-b_n}\\
	b_{n+1}&=&\text{e}^{-a_n}
\end{array}\right.\]

\begin{enumerate}
	\item 
	\begin{enumerate}
		\item $a_1=\e^{-b_0}=\e^{-1}\approx 0,37$ \quad et \quad $b_1=\e^{a_0}=\e^{-\frac{1}{10}}\approx 0,90$.
		\item Pour utiliser en Python la fonction exponentielle, il faut charger en premier la fonction \textit{exp()} de la librairie "math" : \textit{from math import exp}. 
		
		\begin{center}
			
			\begin{tabular}{|l|}\hline
				\texttt{from math import exp}\\
				\texttt{def termes(n) :}\\
				\quad \texttt{a = 1/10}\\
				\quad \texttt{b = 1}\\
				\quad \texttt{for k in range(0, n) :}\\
				\qquad \texttt{c =  exp(-b)}\\
				\qquad \texttt{b  = exp(-a)}\\
				\qquad \texttt{a  =  c}\\
				\quad \texttt{return(a, b)}\\ \hline
			\end{tabular}
		\end{center}
	\end{enumerate}
\item  

\begin{enumerate}
	\item Montrons par récurrence que $\forall n \in \N$, $0 < a_n \leqslant a_{n+1} \leqslant b_{n+1} \leqslant b_n \leqslant 1$.

\textit{Initialisation :} $a_0=\dfrac{1}{10}$ \quad $a_1=\e^{-1}=\dfrac{1}{\e}$ \quad $b_0=1=\e^0$  \quad $b_1=\e^{-\frac{1}{10}}$.

La fonction $x\mapsto\e^x$ est strictement croissante sur $\R$ et  $-1\leqslant -\dfrac{1}{10} \leqslant 0$ donc ,  $\e^{-1}\leqslant \e^{-\frac{1}{10}} \leqslant \e^0$ et $\dfrac{1}{\e}>\dfrac{1}{10}$ donc  on peut alors affirmer que $0<a_0 \leqslant a_1 \leqslant b_1 \leqslant b_0 \leqslant1$.

L'initialisation est vérifiée.

\textit{Hérédité :} Soit  $n \in \N^*$, et supposons que $0 < a_n \leqslant a_{n+1} \leqslant b_{n+1} \leqslant b_n \leqslant 1$.
		
Montrons que $0 < a_{n+1} \leqslant a_{n+2} \leqslant b_{n+2} \leqslant b_{n+1} \leqslant 1$.
					
La fonction $x \longmapsto \text{e}^{-x}$ est décroissante sur $\R$. Donc 
			
$0 < a_n \leqslant a_{n+1} \leqslant b_{n+1} \leqslant b_n \leqslant 1 \iff \e^{-0} >\e^{-a_n }\geqslant \e^{-a_{n+1}}\geqslant \e^{-b_{n+1}} \geqslant \e^{-b_n} \geqslant \e^{-1} $
			
soit $0< \e^{-1} \leqslant \e^{-b_n} \leqslant \e^{-b_{n+1}} \leqslant \e^{-a_{n+1}} \leqslant \e^{-a_n } < 1 \leqslant 1$ donc

$0 < a_{n+1} \leqslant a_{n+2} \leqslant b_{n+2} \leqslant b_{n+1} \leqslant 1$. On obtient ce qu'il fallait démontrer.
			
L'hérédité est démontrée.
					
\textit{Conclusion : } La proposition est vraie au rang $0$ et si elle est vraie au rang $n$, elle l'est aussi au rang $n + 1$.	D'après l'axiome de récurrence, $\forall n \in \N$, $0 < a_n \leqslant a_{n+1} \leqslant b_{n+1} \leqslant b_n \leqslant 1$.
			\item Nous venons de montrer que $\forall n \in \N$, $0<a_n\leqslant a_{n+1} \leqslant1$ et que $0 \leqslant b_{n+1} \leqslant b_n  \leqslant 1$.
			
			La suite $(a_n)$ est croissante et majorée par 1. D'après  le théorème de convergence monotone, la suite $(a_n)$ converge.
			
			De même, la suite $(b_n)$ est décroissante et minorée par 0. D'après  le théorème de convergence monotone, la suite $(b_n)$ converge.
\end{enumerate}
\item 
	\begin{enumerate}
		\item $A = \text{e}^{-B}$ et $B = \text{e}^{-A}$.

		$A = \text{e}^{-B} \iff \ln(A)=-B=-\e^{-A}$ donc $\ln(A)+\e^{-A}=0$ donc $f(A)=0$.

		\item De même $B = \e^{-A}\iff \ln(B)=-A=-\e^{B}$ donc $\ln(B)+\e^{-B}=0$ donc $f(B)=0$.

Donc A et B sont deux solutions de l'équation $f(x)=0$. Or d'après la question A.4, cette équation admet une unique solution, donc $A= B$ et $A - B = 0$.
	\end{enumerate}
\end{enumerate}

\end{document}
